% TOP_PROBLEM: {"id": 2100202, "title": "Open Problems in Integrable Systems — Topology of integrable systems, Lagrangian fibrations, and their invariants", "category_id": 11, "category": {"id": 11, "name": "dynamical_systems", "display_name": "Dynamical Systems", "description": "Problems about long-term behavior of deterministic systems, Hamiltonian dynamics, and periodic orbits.", "slug": "dynamical-systems", "order_index": 11, "created_at": "2026-07-31T15:26:25.671Z"}, "status": "partially_solved"}
% FIRST_POSTED: null
% ENTRY_KIND: statement_only
% Imported from: batch08.tex; UnsolvedMath ID 2100202
% Compile twice: lualatex 2100202.tex
\documentclass[11pt]{article}
\usepackage{fontspec}
\setmainfont{Latin Modern Roman}
\usepackage{amsmath,amssymb,mathtools,unicode-math}
\setmathfont{Latin Modern Math}
\usepackage{xurl,hyperref}
\hypersetup{hidelinks,pdftitle={UnsolvedMath Problem 2100202}}
\newcommand{\sref}[2]{\hyperlink{source:#1:#2}{[S#2]}}
\newcommand{\cataloguescope}[1]{\paragraph{Mathematical question.}#1}
\begin{document}
% UnsolvedMath ID 2100202; source catalogue code AMR-020-0202
\setcounter{section}{2100201}
\section{Open Problems in Integrable Systems — Topology of integrable systems, Lagrangian fibrations, and their invariants}\label{2100202}
{\small\sffamily UnsolvedMath ID 2100202 · Dynamical Systems}
\subsection{Definitions and mathematical statement}

\paragraph{Shared notation.}$\mathbb N=\{1,2,\ldots\}$, and $\mathbb Z,\mathbb Q,\mathbb R,\mathbb C$ denote the integers, rationals, reals and complex numbers. Any other conventions are specified locally.


A smooth integrable system on a symplectic $2n$-manifold $(M,\omega)$ is a map $F=(F_1,\ldots,F_n)$ whose components Poisson commute and have linearly independent differentials at almost every point, with respect to the symplectic volume measure $\omega^n/n!$. The Poisson bracket is $\{u,v\}=\omega(X_u,X_v)$, where $\iota_{X_u}\omega=-du$. The connected components of common level sets form a singular Lagrangian fibration. On a compact regular fiber, choose a cycle basis $\gamma_i$ and a local primitive $d\alpha=\omega$; the action coordinates are $I_i=(2\pi)^{-1}\oint_{\gamma_i}\alpha$, with changes in $GL(n,\mathbb Z)\ltimes\mathbb R^n$. They define the integral affine structure of the regular base. Nondegeneracy means that the transverse quadratic forms of the integrals generate a Cartan subalgebra of the symplectic linear Lie algebra. Their Williamson normal form is a direct sum of elliptic blocks $(q^2+p^2)/2$, hyperbolic blocks $qp$, and focus--focus pairs $(q_1p_2-q_2p_1,q_1p_1+q_2p_2)$. In dimension four an almost toric fibration has only nondegenerate elliptic and focus--focus singularities. The source considers a two-dimensional base domain with boundary and corners, with isolated focus singularities, carrying the action-induced affine structure. An affine equivalence of such bases is a stratified diffeomorphism preserving this singular affine structure; on regular charts it is an integral affine map.
\paragraph{Statement recorded in the supplied source.}
For two such almost toric fibrations $\phi:(M^4,\omega)\to B$ and $\phi':({M'}^4,\omega')\to B'$ with compact fibers, does an affine diffeomorphism $B\to B'$ imply that the fibrations are symplectomorphic? Here equivalence of fibrations means a symplectomorphism sending fibers to fibers. \sref{2100202}{2}
\subsection{Short English statement}
Does the singular integral affine base determine an almost toric four-dimensional Lagrangian fibration up to symplectic equivalence?
\subsection{Sources}
\begin{description}
\item[\hypertarget{source:2100202:1}{S1}] ULAM.AI and the UnsolvedMath Contributors, UnsolvedMath: A Curated Collection of Open Mathematics Problems, record 2100202 (AMR-020-0202). CC BY 4.0. \url{https://huggingface.co/datasets/ulamai/UnsolvedMath}.
\item[\hypertarget{source:2100202:2}{S2}] A. Bolsinov, V. S. Matveev, E. Miranda and S. Tabachnikov, Open Problems, Questions, and Challenges in Finite-Dimensional Integrable Systems (2018), arXiv:1804.03737, Problem 2.2 and preceding almost-toric paragraph. \url{https://arxiv.org/abs/1804.03737}.
\end{description}
\label{end:2100202}
\end{document}
